% Abhakseite „Das kann ich“ – Zone nur im Gesamt (\mitzone)
%% TODO Ich-kann-Sätze: Platzhalter wie in den Aufgabendateien
\begin{abhakseite}
\ifdefined\mitzone
\abhakgruppe{Kennst du schon}
\abhak{1}{Teilen und Vervielfachen im Kopf (vierundzwanzig geteilt durch sechs, drei mal neun), Einmaleins}
\abhak{2}{Größen mit Komma lesen und umrechnen (1,2 kg = 1200 g; 2,50 €)}
\abhak{3}{Kreissektor als Anteil vom Vollkreis (120° von 360°)}
\abhak{4}{Teiler und Vielfache einer Zahl angeben (Teiler von zwölf; Vielfache von vier)}
\abhak{5}{Stellenwerte natürlicher Zahlen und Zahlenstrahl mit natürlichen Zahlen (Skala ablesen)}
\abhak{6}{Schriftlich oder mit Taschenrechner dividieren (drei geteilt durch acht)}
\abhak{7}{Prozent als Hundertstel (45 \% = 0,45)}
\abhak{8}{Prozent als Hundertstel (45 \% = 0,45) – Fehler finden}
\abhak{9}{Prozent als Hundertstel (45 \% = 0,45) – selbst rechnen}
\fi
\abhakgruppe{Einheit 1 · Bruch als Anteil}
\abhak{10}{Anteil bestimmen}
\abhak{11}{Bruchteil einer Größe}
\abhak{12}{Ganzes aus Bruchteil}
\abhak{13}{Anteil bestimmen – Fehler finden, Begründen, Darstellungswechsel, Anwendung}
\abhakgruppe{Einheit 2 · Kürzen und Erweitern}
\abhak{14}{Kürzen und Erweitern}
\abhak{15}{gleichwertige Brüche am Streifen finden}
\abhak{16}{Bruch als Geteilt-Aufgabe}
\abhak{17}{Kürzen und Erweitern – Fehler finden, Begründen, Darstellungswechsel, Anwendung}
\abhakgruppe{Einheit 3 · Brüche vergleichen}
\abhak{18}{Vergleichen}
\abhak{19}{Zahl zwischen zwei Brüchen angeben}
\abhak{20}{Vergleichen – Fehler finden, Begründen, Darstellungswechsel, Anwendung}
\abhakgruppe{Einheit 4 · Dezimalzahlen}
\abhak{21}{Umwandeln}
\abhak{22}{Zahl in die Stellenwerttafel eintragen und ablesen (Zehntel, Hundertstel, Tausendstel)}
\abhak{23}{Nachbarzahlen (Nachbar-Einer, -Zehntel, -Hundertstel) und Zählen in Schritten (0,2er-Schritte; über 2,9 hinaus)}
\abhak{24}{Umwandeln – Fehler finden, Begründen, Darstellungswechsel, Anwendung}
\abhakgruppe{Einheit 5 · Vergleichen, Ordnen, Runden}
\abhak{25}{Welche Form?}
\abhak{26}{Vergleichen und Ordnen}
\abhak{27}{Zahl zwischen zwei Zahlen angeben (auch zwischen zwei Brüchen über Dezimalzahlen)}
\abhak{28}{gemischte Liste ordnen (Bruch, Dezimalzahl, Prozent, Potenz einer Dezimalzahl)}
\abhak{29}{Aussagen prüfen und die wahre ankreuzen}
\abhak{30}{Vergleichen und Ordnen – Fehler finden, Begründen, Darstellungswechsel, Anwendung}
\end{abhakseite}
